% ============================================================
%  Rapport de projet — Analyse des collisions routières à Montréal
%  [DRAFT À COMPLÉTER]
% ============================================================
\documentclass[12pt]{article}
 
% --- Encodage et langue ---
\usepackage[T1]{fontenc}
\usepackage[utf8]{inputenc}
\usepackage[french]{babel}
 
% --- Mise en page ---
\usepackage[margin=1in]{geometry}
\usepackage{setspace}
\onehalfspacing          % interligne 1.5 — ajuster si besoin
 
% --- Figures, tableaux ---
\usepackage{graphicx}
\usepackage{booktabs}
\usepackage{float}
\usepackage{caption}
\usepackage{subcaption}
 
% --- Mathématiques ---
\usepackage{amsmath}
\usepackage{amssymb}
 
% --- Liens hypertexte ---
\usepackage[hidelinks]{hyperref}
 
% --- Code source ---
\usepackage{listings}
\usepackage{xcolor}
\lstset{
  basicstyle=\small\ttfamily,
  breaklines=true,
  frame=single,
  backgroundcolor=\color{gray!10},
  keywordstyle=\color{blue},
  commentstyle=\color{green!50!black},
  stringstyle=\color{orange},
}
 
% --- Bibliographie ---
\usepackage[backend=biber, style=numeric, sorting=none]{biblatex}
\addbibresource{references.bib}   % créer ce fichier avec vos références
 
% ============================================================
\begin{document}
 
% ============================================================
%  Page de titre
% ============================================================
\begin{titlepage}
  \centering
  \vspace*{2cm}
  {\LARGE \textbf{Analyse, regroupement et classification\\
  des collisions routières à Montréal}\\[0.5em]}
  \rule{\linewidth}{0.4pt}\\[1.5em]
  {\large Rapport de projet}\\[1em]
  {\large Fondements en sciences des données}\\[2em]
  {\large  STT3795}\\[2em]

  {\large
  Xavier Drolet (20276341)\\
  Benjamin Dockins (20267474)\\
  Pierre Emery (20278920)\\[2em]}
 
  {\large 23 avril 2026}
  \vfill
  {\large Université de Montréal}
\end{titlepage}
 
% ============================================================
%  Table des matières
% ============================================================
\tableofcontents
\newpage
 
% ============================================================
\section{Introduction}
% ============================================================
 
% TODO : Rédiger l'introduction (contexte général de la sécurité
% routière à Montréal, importance du jeu de données ouvert, etc.)
%
% Pistes à développer :
%  - Contexte sociétal : accidents de la route, blessés, décès à Mtl
%  - Disponibilité des données ouvertes de la Ville de Montréal
%  - Pertinence de l'apprentissage automatique pour ce domaine
 
La sécurité routière constitue un enjeu majeur pour les villes modernes.
À Montréal, les collisions routières occasionnent chaque année de nombreux
blessés et des dommages matériels considérables.
% [COMPLÉTER : statistiques récentes sur les accidents à Montréal]
 
La Ville de Montréal met à disposition un jeu de données ouvert recensant
l'ensemble des collisions signalées sur son territoire depuis 2012~\cite{donnees_mtl}.
Ce jeu de données, riche en variables spatiales, temporelles et contextuelles,
offre un terrain particulièrement propice à l'application de méthodes
d'apprentissage automatique supervisées et non supervisées.
 
% [COMPLÉTER : transition vers la présentation du rapport]
 
 
% ============================================================
\section{Objectifs}
% ============================================================
 
Il y a deux objectifs qui sont poursuivis dans ce projet, en premier lieu on veut identifier différents profils de collisions à Montréal en regroupant les
accidents selon leurs caractéristiques communes : contexte spatial, moment de survenue, conditions environnementales et type d'accident.
Il s'agit d'évaluer dans quelle mesure des méthodes de regroupement permettent de faire ressortir des structures naturelles dans les données. D'autre part, il y aussi un travail de classification supervisé, le but est de prédire la gravité d'une collision à partir de ses attributs.
Nous souhaitons comparer plusieurs approches afin de déterminer lesquelles
sont les mieux adaptées à ce type de données.
 
%\textbf{Objectif 3 — Prédiction spatiale.}
%Si le temps le permet, modéliser la probabilité de collision grave selon
%la localisation géographique afin d'identifier les zones les plus
%dangereuses du réseau routier montréalais.
 
 
% ============================================================
\section{Description des données}
% ============================================================
 
\subsection{Source et structures des données}
 
Les données proviennent du jeu de données ouvert de la Ville de Montréal sur les collisions routières~\cite{donnees_mtl}, accessible à
l'adresse \url{https://donnees.montreal.ca/dataset/collisions-routieres}. 
% Un second jeu de données portant sur le comptage de véhicules, piétons et
% cyclistes~\cite{comptage_mtl} pourrait également être intégré pour enrichir
% certaines analyses. 
 
Le jeu de données brut compte 218\,272 entrées et 68 colonnes couvrant la période 2012--2021.
Chaque entrée correspond à une collision individuelle. Les variables se répartissent en plusieurs catégories :
 
\begin{itemize}
  \item {Variables temporelles :} date (\texttt{DT\_ACCDN}),
        heure (\texttt{HEURE\_ACCDN}), jour de la semaine
        (\texttt{JR\_SEMN\_ACCDN}), année (\texttt{AN}).
  \item {Variables spatiales :} coordonnées géographiques
        (\texttt{LOC\_LAT}, \texttt{LOC\_LONG}), numéro civique,
        nom de rue, région administrative, code municipal.
  \item {Variables contextuelles et environnementales :}
        conditions météorologiques (\texttt{CD\_COND\_METEO}),
        état de la chaussée (\texttt{CD\_ETAT\_SURFC}),
        éclairage (\texttt{CD\_ECLRM}),
        configuration de la route (\texttt{CD\_CONFG\_ROUTE}),
        vitesse autorisée (\texttt{VITESSE\_AUTOR}), etc.
  \item {Variable cible :} gravité de la collision
        (\texttt{GRAVITE}), avec cinq niveaux d'origine.
\end{itemize}
 
\subsection{Prétraitement et nettoyage des données}
 
Le nettoyage des données a été réalisé en plusieurs étapes successives,
implémentées dans le module \texttt{data\_cleaning.py}.
 
\subsubsection{Suppression de colonnes}
 
Trois catégories de colonnes ont été retirées du jeu de données.
 
\paragraph{Colonnes avec trop de valeurs manquantes.}
Dix colonnes présentant plus de 60\,\% de valeurs manquantes ont été
supprimées, car leur imputation aurait introduit trop d'incertitude :
\texttt{BORNE\_KM\_ACCDN}, \texttt{CD\_PNT\_CDRNL\_ROUTE},
\texttt{SFX\_NO\_CIVIQ\_ACCDN}, \texttt{NO\_ROUTE},
\texttt{CD\_ZON\_TRAVX\_ROUTR}, \texttt{CD\_SIT\_PRTCE\_ACCDN},
\texttt{CD\_POSI\_ACCDN}, \texttt{CD\_ETAT\_CHASS},
\texttt{NB\_METRE\_DIST\_ACCD} et \texttt{CD\_PNT\_CDRNL\_REPRR}.
 
\paragraph{Colonnes d'identification et textuelles.}
Six colonnes jugées sans valeur analytique ont également été retirées :
l'identifiant unique de collision (\texttt{NO\_SEQ\_COLL}),
le code municipal redondant avec la région administrative
(\texttt{CD\_MUNCP}), le numéro civique trop granulaire
(\texttt{NO\_CIVIQ\_ACCDN}, 57\,\% de valeurs manquantes),
le nom de rue (\texttt{RUE\_ACCDN}, trop de valeurs uniques),
l'intersection proche (\texttt{ACCDN\_PRES\_DE}, texte libre,
33\,\% de NaN) et le type de repère (\texttt{TP\_REPRR\_ACCDN},
38\,\% de NaN, peu informatif).
 
\paragraph{Colonnes redondantes.}
Les colonnes \texttt{LOC\_X} et \texttt{LOC\_Y} ont été supprimées
car elles sont redondantes avec \texttt{LOC\_LONG} et \texttt{LOC\_LAT}.
 
\subsubsection{Traitement des variables temporelles}
 
La colonne \texttt{DT\_ACCDN} (format \texttt{YYYY/MM/DD}) a été
transformée en deux nouvelles variables : le mois (\texttt{MOIS},
entier de 1 à 12) et la saison (\texttt{SAISON} : Hiver, Printemps,
Été, Automne). La colonne \texttt{HEURE\_ACCDN}, encodée sous forme
de plages horaires (ex. \texttt{"02:00:00-02:59:00"}) ou de la valeur
\texttt{"Non précisé"}, a été convertie en entier représentant
l'heure de début (0--23). Les colonnes sources ont ensuite été supprimées.
 
\subsubsection{Retypage des codes catégoriels}
 
Neuf colonnes de type \texttt{CD\_*} sont encodées sous forme de codes
numériques représentant en réalité des catégories (genre d'accident,
état de la surface, éclairage, etc.). Ces colonnes ont été converties
de \texttt{float64} vers le type \texttt{string} afin d'éviter que les
algorithmes d'apprentissage automatique ne les traitent comme des
variables continues.
 
\subsubsection{Imputation des valeurs manquantes}
 
La stratégie d'imputation est adaptée à chaque variable.
 
\paragraph{Coordonnées géographiques.}
Puiqu'il y a très peu de lignes présentant des valeurs manquantes dans
\texttt{LOC\_LONG} ou \texttt{LOC\_LAT}, elles ont simplement été supprimées, puisque l'effet sur le jeu de données est minime.
 
\paragraph{Vitesse autorisée (\texttt{VITESSE\_AUTOR}) et l'heure (\texttt{HEURE}).}
Dans ces deux cas, il y avait une portion considérables des valeurs de ces colonnes qui étaient manquante (37\% et ? \%). Nous pensions d'abord faire l'imputation par la médiane globale ou bien par le mode mais cette stratégie créait un pic artificiel de taille démesurée
non représentatif de la distribution réelle. Nous avons donc opté
pour un tirage aléatoire pondéré selon la distribution de la variable. Donc pour la vitesse autorisée, 
au sein de chaque catégorie de route (\texttt{CD\_CATEG\_ROUTE}),
avec un générateur aléatoire à graine fixe (\texttt{seed=3795}) pour
la reproductibilité. En cas de catégorie de route elle-même manquante,
la distribution globale des valeurs connues a servi de recours.
Puis, les heures manquantes (valeurs \texttt{"Non précisé"}
dans les données brutes) ont été imputées par tirage aléatoire selon la
distribution empirique des heures connues.
 
\paragraph{Variables catégorielles restantes.}
Pour toutes les colonnes de type \texttt{object} ou \texttt{string}
présentant encore des valeurs manquantes après les étapes précédentes,
les \texttt{NaN} ont été remplacés par la catégorie \texttt{"Inconnu"},
qui est ainsi traitée comme une modalité à part entière.
 
\subsubsection{Suppression des doublons}
 
Les lignes entièrement dupliquées ont été supprimées.
 
\subsubsection{Regroupement de la variable cible}
 
La variable \texttt{GRAVITE} comporte originellement cinq niveaux.
Ces niveaux ont été regroupés en trois classes pour la tâche de
classification supervisée :
 
\begin{itemize}
  \item \textbf{Grave} : regroupe \emph{Mortel} et
        \emph{Blessures graves}.
  \item \textbf{Léger} : correspond à \emph{Blessures légères}.
  \item \textbf{Matériel} : regroupe \emph{Dommages matériels seulement}
        et \emph{Dommages matériels inférieurs au seuil de rapportage}.
\end{itemize}
 
Ce regroupement se justifie par la faible représentation des accidents
mortels et graves, qui risquerait de nuire à l'apprentissage d'un
classifieur, et par la similarité fonctionnelle des deux sous-catégories
de dommages matériels.
 
% ============================================================
\section{Méthodologie}
% ============================================================

 
% TODO : Ajouter le tableau de distribution des classes après regroupement
 
\subsection{Analyse exploratoire}
 
% TODO : Décrire brièvement les visualisations produites dans
% data_exploration.ipynb (distribution des variables, carte des accidents,
% corrélations, etc.)
 
\subsection{Réduction de dimension}
 
% TODO : Décrire la PCA et éventuellement MDS / Isomap
% Notions du cours à mentionner : valeurs propres, variance expliquée,
% biplot, etc.
 
\subsection{Regroupement (clustering)}
 
% TODO : Décrire les méthodes utilisées (k-moyennes, clustering
% hiérarchique) et les critères de sélection du nombre de groupes
% (coude, silhouette, dendrogramme)
 
\subsection{Classification supervisée}
 
% TODO : Décrire les classifieurs comparés :
%   - Naive Bayes (baseline)
%   - SVM
%   - Arbres de décision / Random Forest
% Mentionner la validation croisée, les métriques (accuracy, F1, etc.)
% et le déséquilibre des classes
 
 
% ============================================================
\section{Résultats}
% ============================================================
 
% TODO : Présenter les résultats de chaque méthode.
% Cette section peut inclure des tentatives infructueuses.
 
\subsection{Nettoyage et exploration}
 
% TODO : Résumé quantitatif du nettoyage (nb lignes/colonnes avant/après,
% distribution des classes GRAVITE_3, etc.)
 
\subsection{Réduction de dimension}
 
% TODO : Variance expliquée par les composantes principales,
% visualisations 2D/3D
 
\subsection{Regroupement}
 
% TODO : Résultats des k-moyennes et du clustering hiérarchique,
% interprétation des groupes identifiés
 
\subsection{Classification supervisée}
 
% TODO : Tableau comparatif des performances des classifieurs
% (accuracy, F1-macro, matrice de confusion, etc.)
 
 
% ============================================================
\section{Conclusion et discussion}
% ============================================================
 
% TODO : Synthétiser ce qui a fonctionné et ce qui n'a pas fonctionné.
% Discuter des limites (déséquilibre des classes, qualité des données,
% variables manquantes). Proposer des directions futures.
 
% Pistes possibles :
%  - Intégration du jeu de données de comptage de véhicules/piétons
%  - Modèle de prédiction spatiale (objectif 3)
%  - Features supplémentaires (météo externe, densité de population)
%  - Méthodes d'ensemble (XGBoost, Random Forest avec tuning)
 
 
% ============================================================
\section{Contributions des membres de l'équipe}
% ============================================================
 
\paragraph{Xavier Drolet (20276341).}
% TODO : Décrire les contributions de Xavier
% (selon la proposition : collecte/organisation des données,
%  puis méthodes de réduction de dimension)
 
\paragraph{Benjamin Dockins (20267474).}
% TODO : Décrire les contributions de Benjamin
% (selon la proposition : nettoyage/prétraitement des données,
%  puis méthodes de classification supervisée)
 
\paragraph{Pierre Emery (20278920).}
% TODO : Décrire les contributions de Pierre
% (selon la proposition : nettoyage et extraction de features,
%  puis méthodes de regroupement)
 
 
% ============================================================
%  Bibliographie
% ============================================================
\printbibliography
 
\end{document}
 
% ============================================================
%  Exemple de fichier references.bib à créer :
% ============================================================
%
% @misc{donnees_mtl,
%   author       = {{Ville de Montréal}},
%   title        = {Collisions routières},
%   year         = {2024},
%   howpublished = {\url{https://donnees.montreal.ca/dataset/collisions-routieres}},
%   note         = {Consulté le [DATE]}
% }
%
% @misc{comptage_mtl,
%   author       = {{Ville de Montréal}},
%   title        = {Comptage de véhicules, piétons et cyclistes},
%   year         = {2024},
%   howpublished = {\url{https://donnees.montreal.ca/dataset/comptage-vehicules-pietons}},
%   note         = {Consulté le [DATE]}
% }
 
